\emph{ट्यूरिंग यंत्र} ही आदर्शीकृत संगणन यंत्रणा आहे. तिच्यात एका दिशेने अनंत असलेली \emph{फीत} असते. फीत घरांमध्ये विभागलेली असते आणि प्रत्येक घरात आधी ठरवलेल्या वर्णमालेतील एक \emph{चिन्ह} असू शकते. फितीवर \emph{वाचन-लेखन शीर्ष} हलवून यंत्र काम करते. यंत्राच्या \emph{अवस्थांची} संख्या आधीच ठरलेली असते. यंत्र त्यांपैकी एका अवस्थेत असू शकते. वाचन-लेखन शीर्षाच्या क्रिया सूचनांचा संच ठरवतो. प्रत्येक सूचना यंत्र एका विशिष्ट अवस्थेत असणे आणि त्याने एक विशिष्ट चिन्ह वाचणे यांवर अवलंबून असते. सध्याच्या घरात कोणते चिन्ह लिहायचे, वाचन-लेखन शीर्ष एक घर डावीकडे किंवा उजवीकडे हलवायचे की जागेवरच ठेवायचे आणि यंत्राने कोणत्या अवस्थेत जायचे हे ती सूचना ठरवते. फितीवर चिन्हांच्या क्रमिकेच्या रूपात एखादे \emph{आदान} दिले असेल, यंत्र ठरवून दिलेल्या आरंभीच्या अवस्थेत ठेवले असेल आणि वाचन-लेखन शीर्ष फितीच्या डाव्या टोकाच्या चिन्हालगत उजवीकडील घर वाचत असेल, तर सूचनासंच यंत्राच्या \emph{स्थितिवर्णनांची} क्रमिका पायरीपायरीने ठरवतो. स्थितिवर्णनात फितीवरील मजकूर, वाचन-लेखन शीर्षाचे स्थान आणि यंत्राची अवस्था येतात. एखाद्या स्थितिवर्णनातील वाचलेले चिन्ह आणि सध्याची अवस्था यांच्या जोडीसाठी सूचनासंचात सूचना नसेल, तर यंत्र \emph{थांबते}. डाव्या टोकाचे चिन्ह शाबूत असल्यास ते चिन्ह आणि उजवीकडील शेवटची सलग रिकामी घरे वगळून उरलेली चिन्हमाला हे यंत्राचे \emph{प्रदान} असते.

ट्यूरिंग यंत्राच्या फितीवर संख्यांचे रेघांच्या क्रमिकेच्या रूपात निरूपण अगदी सहज करता येते. फलन $f\colon\Nat \to \Nat$ हे \emph{ट्यूरिंग-संगणनक्षम} असते, जर असे ट्यूरिंग यंत्र अस्तित्वात असेल की जेव्हा जेव्हा ते~$n$ चे एकचिन्ही निरूपण आदान म्हणून घेऊन सुरू होते, तेव्हा ते अखेरीस थांबते आणि प्रदान म्हणून त्याच्या फितीवर~$f(n)$ चे एकचिन्ही निरूपण असते. अंकगणितातील अनेक परिचित फलने ट्यूरिंग-संगणनक्षम आहेत हे दाखवता येते, काहींसाठी ते सहज जमते तर काहींसाठी अधिक प्रयत्न लागतात. ट्यूरिंग यंत्रांखेरीज संगणनाची इतरही अनेक प्रतिमाने सुचवली गेली आहेत. त्या प्रतिमानांमध्ये संगणनक्षम असलेली अंकगणितीय फलने ट्यूरिंग-संगणनक्षमही असतात, असे प्रत्येक वेळी दिसून आले आहे. प्रभावीपणे संगणन करता येणारे प्रत्येक अंकगणितीय फलन ट्यूरिंग-संगणनक्षम आहे, या \emph{चर्च--ट्यूरिंग प्रबंधाला} यातून आधार मिळतो.
